% ============================================================================
% CAPÍTULO 9: SIGNALING SYSTEMS
% Bioinformática para Biologia de Sistemas
% ============================================================================

% ============================================================================
% TEMPLATE BASE PARA APRESENTAÇÕES EM BEAMER
% Bioinformática para Biologia de Sistemas
% ============================================================================

\documentclass[12pt, aspectratio=169]{beamer}

% ============================================================================
% PACOTES ESSENCIAIS
% ============================================================================
\usepackage[utf8]{inputenc}
\usepackage[T1]{fontenc}
\usepackage[brazilian]{babel}
\usepackage{amsmath, amssymb, amsthm}
\usepackage{graphicx}
\usepackage{xcolor}
\usepackage{tikz}
\usetikzlibrary{arrows.meta, shapes, positioning, calc, decorations.pathreplacing, decorations.shapes, decorations.pathmorphing}
\usepackage{listings}
\usepackage{booktabs}
\usepackage{multirow}
\usepackage{hyperref}
\usepackage{chemfig}
\usepackage{chemarr}

% ============================================================================
% CONFIGURAÇÃO DO TEMA
% ============================================================================
\usetheme{Madrid}
\usecolortheme{default}

% Cores personalizadas
\definecolor{azulescuro}{RGB}{0, 51, 102}
\definecolor{azulclaro}{RGB}{51, 153, 255}
\definecolor{verdeacido}{RGB}{102, 204, 0}
\definecolor{laranjacel}{RGB}{255, 153, 51}
\definecolor{roxodna}{RGB}{153, 51, 255}

\setbeamercolor{structure}{fg=azulescuro}
\setbeamercolor{palette primary}{bg=azulescuro,fg=white}
\setbeamercolor{palette secondary}{bg=azulclaro,fg=white}
\setbeamercolor{palette tertiary}{bg=azulescuro,fg=white}
\setbeamercolor{palette quaternary}{bg=azulclaro,fg=white}
\setbeamercolor{block title}{bg=azulescuro,fg=white}
\setbeamercolor{block body}{bg=azulclaro!10,fg=black}
\setbeamercolor{block title alerted}{bg=red!80,fg=white}
\setbeamercolor{block body alerted}{bg=red!10,fg=black}
\setbeamercolor{block title example}{bg=verdeacido!80,fg=white}
\setbeamercolor{block body example}{bg=verdeacido!10,fg=black}

% ============================================================================
% CONFIGURAÇÃO DE CÓDIGO
% ============================================================================
\lstset{
    language=Python,
    basicstyle=\ttfamily\small,
    keywordstyle=\color{azulescuro}\bfseries,
    commentstyle=\color{verdeacido},
    stringstyle=\color{laranjacel},
    numbers=left,
    numberstyle=\tiny\color{gray},
    stepnumber=1,
    numbersep=5pt,
    backgroundcolor=\color{white},
    showspaces=false,
    showstringspaces=false,
    showtabs=false,
    frame=single,
    rulecolor=\color{black},
    tabsize=2,
    captionpos=b,
    breaklines=true,
    breakatwhitespace=false,
    escapeinside={\%*}{*)},
    xleftmargin=10pt,
    xrightmargin=10pt
}

% Definir estilo para R
\lstdefinestyle{Rstyle}{
    language=R,
    basicstyle=\ttfamily\small,
    keywordstyle=\color{azulescuro}\bfseries,
    commentstyle=\color{verdeacido},
    stringstyle=\color{laranjacel}
}

% Definir estilo para MATLAB
\lstdefinestyle{MATLABstyle}{
    language=Matlab,
    basicstyle=\ttfamily\small,
    keywordstyle=\color{azulescuro}\bfseries,
    commentstyle=\color{verdeacido},
    stringstyle=\color{laranjacel}
}

% ============================================================================
% COMANDOS PERSONALIZADOS
% ============================================================================

% Comando para equações destacadas
\newcommand{\destaque}[1]{%
    \begin{alertblock}{}
        \begin{center}
            #1
        \end{center}
    \end{alertblock}
}

% Comando para definições
\newcommand{\definicao}[2]{%
    \begin{block}{Definição: #1}
        #2
    \end{block}
}

% Comando para exemplos
\newcommand{\exemplo}[2]{%
    \begin{exampleblock}{Exemplo: #1}
        #2
    \end{exampleblock}
}

% Comando para observações importantes
\newcommand{\importante}[1]{%
    \begin{alertblock}{Importante!}
        #1
    \end{alertblock}
}

% Comandos matemáticos úteis
\newcommand{\R}{\mathbb{R}}
\newcommand{\N}{\mathbb{N}}
\newcommand{\Z}{\mathbb{Z}}
\newcommand{\dif}{\mathrm{d}}
\newcommand{\parcial}[2]{\frac{\partial #1}{\partial #2}}
\newcommand{\derivada}[2]{\frac{\dif #1}{\dif #2}}
\newcommand{\vect}[1]{\boldsymbol{#1}}

% Comandos biológicos
\newcommand{\gene}[1]{\textit{#1}}
\newcommand{\proteina}[1]{\textsc{#1}}
\newcommand{\especie}[1]{\textit{#1}}

% ============================================================================
% CONFIGURAÇÃO DE RODAPÉ
% ============================================================================
\setbeamertemplate{footline}{
  \leavevmode%
  \hbox{%
  \begin{beamercolorbox}[wd=.33\paperwidth,ht=2.25ex,dp=1ex,center]{author in head/foot}%
    \usebeamerfont{author in head/foot}\insertshortauthor
  \end{beamercolorbox}%
  \begin{beamercolorbox}[wd=.34\paperwidth,ht=2.25ex,dp=1ex,center]{title in head/foot}%
    \usebeamerfont{title in head/foot}\insertshorttitle
  \end{beamercolorbox}%
  \begin{beamercolorbox}[wd=.33\paperwidth,ht=2.25ex,dp=1ex,right]{date in head/foot}%
    \usebeamerfont{date in head/foot}\insertshortdate{}\hspace*{2em}
    \insertframenumber{} / \inserttotalframenumber\hspace*{2ex}
  \end{beamercolorbox}}%
  \vskip0pt%
}

% ============================================================================
% CONFIGURAÇÃO DE NAVEGAÇÃO
% ============================================================================
\setbeamertemplate{navigation symbols}{}

% ============================================================================
% CONFIGURAÇÃO DE ITEMIZE/ENUMERATE
% ============================================================================
\setbeamertemplate{itemize items}[circle]
\setbeamertemplate{enumerate items}[default]

% ============================================================================
% FIM DO TEMPLATE
% ============================================================================


% ============================================================================
% INFORMAÇÕES DO DOCUMENTO
% ============================================================================
\title[Cap. 9: Signaling Systems]{Capítulo 9: Signaling Systems}
\subtitle{Bioinformática para Biologia de Sistemas}
\author{Ney Lemke}
\institute{DBF-Unesp}
\date{\today}

\begin{document}

% ============================================================================
% SLIDE DE TÍTULO
% ============================================================================
\begin{frame}
\titlepage
\end{frame}

% ============================================================================
% SUMÁRIO
% ============================================================================
\begin{frame}{Sumário}
\tableofcontents
\end{frame}

% ============================================================================
% SEÇÃO 1: INTRODUÇÃO
% ============================================================================
\section{Introdução}

\begin{frame}{9.1 Introdução aos Sistemas de Sinalização}
\begin{block}{Objetivos de Aprendizagem}
\begin{itemize}
    \item Compreender os princípios básicos da sinalização celular
    \item Conhecer os principais tipos de receptores e vias de sinalização
    \item Entender a dinâmica e propriedades emergentes de redes de sinalização
    \item Modelar matematicamente cascatas de sinalização
    \item Simular sistemas de sinalização usando Python
\end{itemize}
\end{block}

\vspace{0.3cm}

\importante{
A sinalização celular é essencial para:
\begin{itemize}
    \item Coordenação entre células
    \item Resposta a estímulos ambientais
    \item Diferenciação e desenvolvimento
    \item Homeostase e sobrevivência
\end{itemize}
}
\end{frame}

\begin{frame}{Papel da Sinalização Celular}
\definicao{Sinalização Celular}{
Processo pelo qual células detectam, transmitem e respondem a sinais químicos ou físicos, permitindo comunicação intra e intercelular.
}

\vspace{0.3cm}

\begin{center}
\begin{tikzpicture}[node distance=2.5cm, auto,
    box/.style={rectangle, draw, fill=azulclaro!20, text width=2cm, text centered, rounded corners, minimum height=1cm, font=\small}]

    \node[box, fill=laranjacel!30] (signal) {Sinal\\Extracelular};
    \node[box, fill=roxodna!30, right of=signal] (receptor) {Receptor};
    \node[box, fill=verdeacido!30, right of=receptor] (transd) {Transdução};
    \node[box, fill=azulclaro!30, right of=transd] (response) {Resposta};

    \draw[->, ultra thick, azulescuro] (signal) -- (receptor);
    \draw[->, ultra thick, azulescuro] (receptor) -- (transd);
    \draw[->, ultra thick, azulescuro] (transd) -- (response);
\end{tikzpicture}
\end{center}

\vspace{0.3cm}

\begin{exampleblock}{Exemplos de Respostas Celulares}
\begin{itemize}
    \item Alteração na expressão gênica
    \item Modificação do metabolismo
    \item Reorganização do citoesqueleto
    \item Proliferação, diferenciação ou apoptose
\end{itemize}
\end{exampleblock}
\end{frame}

\begin{frame}{Tipos de Sinalização}
\begin{columns}
\column{0.5\textwidth}
\textbf{Classificação por Alcance:}

\begin{itemize}
    \item \textbf{Autócrina:} célula responde ao próprio sinal
    \item \textbf{Parácrina:} sinal atua em células vizinhas (curta distância)
    \item \textbf{Endócrina:} hormônios transportados pelo sangue (longa distância)
    \item \textbf{Justácrina:} contato direto entre células
\end{itemize}

\column{0.5\textwidth}
\begin{center}
\begin{tikzpicture}[scale=0.7]
    % Autocrine
    \draw[fill=azulclaro!30] (0,3) circle (0.5);
    \node at (0,3) {\tiny A};
    \draw[->, thick, roxodna] (0.5,3.3) to [out=60, in=120, looseness=8] (0.5,2.7);
    \node[below, font=\tiny] at (0,2.3) {Autócrina};

    % Paracrine
    \draw[fill=azulclaro!30] (3,3) circle (0.5);
    \draw[fill=verdeacido!30] (5,3) circle (0.5);
    \draw[->, thick, laranjacel] (3.5,3) -- (4.5,3);
    \node[below, font=\tiny] at (4,2.3) {Parácrina};

    % Endocrine
    \draw[fill=azulclaro!30] (0,0.5) circle (0.5);
    \draw[fill=verdeacido!30] (5,0.5) circle (0.5);
    \draw[->, thick, red, dashed] (0.5,0.5) to [out=0, in=180] (4.5,0.5);
    \node[below, font=\tiny] at (2.5,-0.2) {Endócrina};
\end{tikzpicture}
\end{center}
\end{columns}

\vspace{0.3cm}

\begin{exampleblock}{Exemplos}
\begin{itemize}
    \item \textbf{Autócrina:} fatores de crescimento em câncer
    \item \textbf{Parácrina:} neurotransmissores na sinapse
    \item \textbf{Endócrina:} insulina, hormônios tireoidianos
    \item \textbf{Justácrina:} sinalização Notch-Delta
\end{itemize}
\end{exampleblock}
\end{frame}

% ============================================================================
% SEÇÃO 2: RECEPTORES E MECANISMOS
% ============================================================================
\section{Receptores e Mecanismos de Ativação}

\begin{frame}{9.2 Receptores Celulares: Visão Geral}
\begin{block}{Classificação de Receptores}
\begin{enumerate}
    \item \textbf{Receptores acoplados à proteína G (GPCRs)}
    \item \textbf{Receptores tirosina quinase (RTKs)}
    \item \textbf{Receptores de canais iônicos}
    \item \textbf{Receptores nucleares}
    \item \textbf{Receptores de citocinas (JAK-STAT)}
\end{enumerate}
\end{block}

\vspace{0.3cm}

\begin{center}
\begin{tikzpicture}[scale=0.8]
    % Membrane
    \draw[thick, double] (0,0) -- (10,0);
    \node[left, font=\tiny] at (0,0.3) {Extracelular};
    \node[left, font=\tiny] at (0,-0.3) {Intracelular};

    % GPCR
    \draw[thick, azulclaro, fill=azulclaro!20] (1,-0.5) rectangle (2,0.5);
    \node[below, font=\tiny] at (1.5,-0.8) {GPCR};

    % RTK
    \draw[thick, verdeacido, fill=verdeacido!20] (3,-0.5) rectangle (4,0.5);
    \node[below, font=\tiny] at (3.5,-0.8) {RTK};

    % Ion channel
    \draw[thick, laranjacel, fill=laranjacel!20] (5,-0.5) rectangle (6,0.5);
    \node[below, font=\tiny] at (5.5,-0.8) {Canal};

    % Nuclear receptor
    \draw[thick, roxodna, fill=roxodna!20] (7.5,-1) circle (0.3);
    \node[below, font=\tiny] at (7.5,-1.5) {Nuclear};
\end{tikzpicture}
\end{center}
\end{frame}

\begin{frame}{Receptores Acoplados à Proteína G (GPCRs)}
\begin{columns}
\column{0.5\textwidth}
\textbf{Características:}
\begin{itemize}
    \item 7 domínios transmembrana
    \item Maior família de receptores ($\sim$800 em humanos)
    \item Acoplados a proteínas G heterotriméricas
    \item Alvos de $\sim$50\% dos fármacos
\end{itemize}

\vspace{0.3cm}

\textbf{Subtipos de Proteína G:}
\begin{itemize}
    \item \textbf{G$_s$:} estimula adenilil ciclase $\uparrow$ cAMP
    \item \textbf{G$_i$:} inibe adenilil ciclase $\downarrow$ cAMP
    \item \textbf{G$_q$:} ativa fosfolipase C $\uparrow$ IP$_3$/DAG
    \item \textbf{G$_{12/13}$:} regula citoesqueleto
\end{itemize}

\column{0.5\textwidth}
\begin{center}
\begin{tikzpicture}[scale=0.65]
    % Membrane
    \draw[thick, double] (-1,0) -- (4,0);

    % GPCR serpentine
    \foreach \x in {0, 0.4, 0.8, 1.2, 1.6, 2.0, 2.4} {
        \draw[thick, azulclaro, fill=azulclaro!30] (\x,-0.7) rectangle (\x+0.3,0.7);
    }

    % Ligand
    \draw[fill=laranjacel!50] (1,1.5) circle (0.3);
    \node[above, font=\tiny] at (1,1.8) {Ligante};
    \draw[->, dashed] (1,1.2) -- (1,0.7);

    % G-protein heterotrimer
    \draw[fill=verdeacido!50] (1.2,-1.5) circle (0.4);
    \node[font=\scriptsize] at (1.2,-1.5) {$\alpha$};
    \draw[fill=roxodna!50] (2,-1.3) circle (0.25);
    \node[font=\scriptsize] at (2,-1.3) {$\beta$};
    \draw[fill=roxodna!50] (2,-1.7) circle (0.25);
    \node[font=\scriptsize] at (2,-1.7) {$\gamma$};

    % GTP/GDP
    \node[font=\tiny] at (1.2,-2.2) {GDP/GTP};
\end{tikzpicture}
\end{center}

\begin{exampleblock}{Exemplos}
\begin{itemize}
    \item Receptores adrenérgicos
    \item Receptores muscarínicos
    \item Receptores de olfato
    \item Rodopsina (visão)
\end{itemize}
\end{exampleblock}
\end{columns}
\end{frame}

\begin{frame}{Mecanismo de Ativação de GPCRs}
\begin{center}
\begin{tikzpicture}[node distance=2cm, scale=0.85, transform shape]
    % Estado inativo
    \node[draw, rectangle, fill=azulclaro!20, text width=3cm, align=center, minimum height=1.5cm] (inactive) {
        \textbf{Estado Inativo}\\
        GPCR + G$\alpha\beta\gamma$-GDP
    };

    % Ligante se liga
    \node[draw, rectangle, fill=laranjacel!20, text width=3cm, align=center, minimum height=1.5cm, right of=inactive, node distance=4cm] (bound) {
        \textbf{Ligante Ligado}\\
        L-GPCR* + G$\alpha\beta\gamma$-GDP
    };

    % Ativação
    \node[draw, rectangle, fill=verdeacido!20, text width=3cm, align=center, minimum height=1.5cm, below of=bound] (active) {
        \textbf{Estado Ativo}\\
        G$\alpha$-GTP + G$\beta\gamma$\\
        \small(separados)
    };

    % GTP hidrolisado
    \node[draw, rectangle, fill=roxodna!20, text width=3cm, align=center, minimum height=1.5cm, below of=inactive] (hydrolyze) {
        \textbf{Terminação}\\
        G$\alpha$-GDP + G$\beta\gamma$\\
        \small(reassociam)
    };

    \draw[->, thick] (inactive) -- node[above, font=\small] {ligante} (bound);
    \draw[->, thick] (bound) -- node[right, font=\small] {troca GDP/GTP} (active);
    \draw[->, thick] (active) -- node[below, font=\small] {GTPase} (hydrolyze);
    \draw[->, thick] (hydrolyze) -- node[left, font=\small] {reciclagem} (inactive);
\end{tikzpicture}
\end{center}

\importante{
G$\alpha$-GTP e G$\beta\gamma$ ativam efetores downstream independentemente!
}
\end{frame}

\begin{frame}{Receptores Tirosina Quinase (RTKs)}
\begin{columns}
\column{0.5\textwidth}
\textbf{Características:}
\begin{itemize}
    \item Domínio extracelular (ligação)
    \item Domínio transmembrana único
    \item Domínio intracelular (quinase)
    \item Ativados por dimerização
    \item Autofosforilação em tirosinas
\end{itemize}

\vspace{0.3cm}

\textbf{Famílias Importantes:}
\begin{itemize}
    \item \textbf{EGFR:} EGF receptor
    \item \textbf{PDGFR:} PDGF receptor
    \item \textbf{FGFR:} FGF receptor
    \item \textbf{InsR:} Insulin receptor
    \item \textbf{VEGFR:} VEGF receptor
\end{itemize}

\column{0.5\textwidth}
\begin{center}
\begin{tikzpicture}[scale=0.7]
    % Membrane
    \draw[thick, double] (-1,0) -- (5,0);

    % Inactive monomers
    \begin{scope}[xshift=-1cm]
        \draw[thick, azulclaro, fill=azulclaro!30] (0,0.5) ellipse (0.5 and 0.8);
        \draw[thick, azulclaro, fill=azulclaro!30] (0,0) rectangle (0.5,-0.3);
        \draw[thick, azulclaro, fill=azulclaro!30] (0.25,-0.8) circle (0.5);
        \node[font=\tiny] at (0.25,-0.8) {K};
    \end{scope}

    \begin{scope}[xshift=1.5cm]
        \draw[thick, azulclaro, fill=azulclaro!30] (0,0.5) ellipse (0.5 and 0.8);
        \draw[thick, azulclaro, fill=azulclaro!30] (0,0) rectangle (0.5,-0.3);
        \draw[thick, azulclaro, fill=azulclaro!30] (0.25,-0.8) circle (0.5);
        \node[font=\tiny] at (0.25,-0.8) {K};
    \end{scope}

    \node[font=\small] at (0.5,1.8) {Inativo};

    % Arrow
    \draw[->, ultra thick] (2.5,0) -- (3.5,0);

    % Active dimer
    \begin{scope}[xshift=4.5cm]
        % Ligand
        \draw[fill=laranjacel!50] (0,1.5) circle (0.3);

        % RTK dimer
        \draw[thick, verdeacido, fill=verdeacido!30] (-0.5,0.5) ellipse (0.5 and 0.8);
        \draw[thick, verdeacido, fill=verdeacido!30] (0.5,0.5) ellipse (0.5 and 0.8);

        \draw[thick, verdeacido, fill=verdeacido!30] (-0.5,0) rectangle (-0.2,-0.3);
        \draw[thick, verdeacido, fill=verdeacido!30] (0.2,0) rectangle (0.5,-0.3);

        \draw[thick, verdeacido, fill=verdeacido!30] (-0.35,-0.8) circle (0.5);
        \draw[thick, verdeacido, fill=verdeacido!30] (0.35,-0.8) circle (0.5);

        % Phosphorylations
        \foreach \y in {-1.2, -1.4, -1.6} {
            \node[font=\scriptsize, red] at (-0.35,\y) {P};
            \node[font=\scriptsize, red] at (0.35,\y) {P};
        }

        \node[font=\small] at (0,2) {Ativo};
    \end{scope}
\end{tikzpicture}
\end{center}
\end{columns}

\vspace{0.3cm}

\exemplo{EGFR em Câncer}{
Mutações/superexpressão de EGFR estão associadas a diversos cânceres. Inibidores de tirosina quinase (TKIs) como gefitinib são usados em terapia.
}
\end{frame}

\begin{frame}{Receptores de Canais Iônicos e Nucleares}
\begin{columns}
\column{0.5\textwidth}
\begin{block}{Canais Iônicos Ligante-Dependentes}
\begin{itemize}
    \item Resposta ultrarrápida (ms)
    \item Alteração do potencial de membrana
    \item Influxo/efluxo de íons
\end{itemize}

\textbf{Exemplos:}
\begin{itemize}
    \item Receptor nicotínico (ACh)
    \item Receptor GABA$_A$
    \item Receptor glutamato (NMDA)
    \item Receptor 5-HT$_3$ (serotonina)
\end{itemize}
\end{block}

\column{0.5\textwidth}
\begin{block}{Receptores Nucleares}
\begin{itemize}
    \item Fatores de transcrição ativados por ligantes
    \item Ligantes hidrofóbicos (atravessam membrana)
    \item Resposta lenta (horas)
\end{itemize}

\textbf{Exemplos:}
\begin{itemize}
    \item Receptores de esteroides (GR, ER, AR)
    \item Receptores de hormônios tireoidianos
    \item Receptores de retinoides (RAR)
    \item PPAR (metabolismo lipídico)
\end{itemize}
\end{block}
\end{columns}

\vspace{0.3cm}

\begin{exampleblock}{Mecanismo de Receptores Nucleares}
Hormônio $\rightarrow$ Atravessa membrana $\rightarrow$ Liga ao receptor no citoplasma/núcleo $\rightarrow$ Dimerização $\rightarrow$ Liga ao DNA (HRE) $\rightarrow$ Ativa/reprime transcrição
\end{exampleblock}
\end{frame}

\begin{frame}{Ligação Ligante-Receptor: Cinética}
\begin{block}{Modelo de Equilíbrio}
\begin{equation*}
L + R \underset{k_{off}}{\overset{k_{on}}{\rightleftharpoons}} LR
\end{equation*}

\textbf{Constante de Dissociação:}
\begin{equation*}
K_d = \frac{k_{off}}{k_{on}} = \frac{[L][R]}{[LR]}
\end{equation*}

\textbf{Fração de Receptores Ocupados:}
\begin{equation*}
\theta = \frac{[LR]}{[R]_{total}} = \frac{[L]}{K_d + [L]}
\end{equation*}
\end{block}

\begin{columns}
\column{0.5\textwidth}
\begin{exampleblock}{Interpretação de $K_d$}
\begin{itemize}
    \item $K_d$ baixo $\rightarrow$ alta afinidade
    \item Quando $[L] = K_d$, $\theta = 0.5$
    \item $K_d$ típicos: nM - $\mu$M
\end{itemize}
\end{exampleblock}

\column{0.5\textwidth}
\begin{exampleblock}{Capacidade Máxima}
\begin{equation*}
B_{max} = [R]_{total}
\end{equation*}
Número máximo de sítios de ligação disponíveis
\end{exampleblock}
\end{columns}
\end{frame}

\begin{frame}[fragile]{Modelagem de Ligação Receptor-Ligante}
\begin{lstlisting}[language=Python, caption=Curva de ligação, basicstyle=\ttfamily\footnotesize]
import numpy as np
import matplotlib.pyplot as plt

def receptor_binding(L, Kd, Bmax):
    """Calcula ligacao pelo modelo de equilibrio"""
    return (Bmax * L) / (Kd + L)

# Parametros
Kd = 10  # nM
Bmax = 100  # unidades arbitrarias

# Concentracoes de ligante
L = np.logspace(-1, 3, 100)  # 0.1 - 1000 nM
bound = receptor_binding(L, Kd, Bmax)

# Plotar
plt.figure(figsize=(10, 6))
plt.semilogx(L, bound, 'b-', linewidth=2)
plt.axhline(y=Bmax/2, color='r', linestyle='--',
            label=f'50% ocupacao')
plt.axvline(x=Kd, color='r', linestyle='--',
            label=f'Kd = {Kd} nM')
plt.xlabel('[Ligante] (nM)')
plt.ylabel('Receptores Ocupados')
plt.title('Curva de Ligacao Ligante-Receptor')
plt.legend()
plt.grid(True, alpha=0.3)
\end{lstlisting}
\end{frame}

\begin{frame}{Cinética de Ativação de Receptores}
\begin{block}{Modelo com Estado Ativo/Inativo}
\begin{equation*}
\begin{aligned}
\frac{d[R]}{dt} &= -k_{on}[L][R] + k_{off}[LR] + k_{inact}[R^*]\\
\frac{d[LR]}{dt} &= k_{on}[L][R] - k_{off}[LR] - k_{act}[LR]\\
\frac{d[R^*]}{dt} &= k_{act}[LR] - k_{inact}[R^*]
\end{aligned}
\end{equation*}

Onde $R^*$ é o receptor ativo que inicia a sinalização.
\end{block}

\begin{exampleblock}{Estado Estacionário}
No equilíbrio ($\frac{d}{dt} = 0$):
\begin{equation*}
[R^*]_{ss} = \frac{k_{act} \cdot k_{on}[L][R]_{total}}{k_{inact}(k_{off} + k_{act}) + k_{act} \cdot k_{on}[L]}
\end{equation*}
\end{exampleblock}
\end{frame}

% ============================================================================
% SEÇÃO 3: REDES DE TRANSDUÇÃO
% ============================================================================
\section{Redes de Transdução de Sinal}

\begin{frame}{9.3 Redes de Transdução: Visão Geral}
\begin{block}{Principais Vias de Sinalização}
\begin{enumerate}
    \item \textbf{MAPK/ERK:} proliferação, diferenciação
    \item \textbf{PI3K/Akt:} sobrevivência, metabolismo
    \item \textbf{JAK-STAT:} resposta a citocinas
    \item \textbf{Wnt:} desenvolvimento, polaridade celular
    \item \textbf{Notch:} decisões de destino celular
    \item \textbf{TGF-$\beta$/Smad:} crescimento, diferenciação
    \item \textbf{NF-$\kappa$B:} inflamação, imunidade
    \item \textbf{Hedgehog:} padronização no desenvolvimento
\end{enumerate}
\end{block}

\importante{
Estas vias não operam isoladamente! Existe extenso \textbf{crosstalk} (diálogo cruzado) entre elas.
}
\end{frame}

\begin{frame}{Via MAPK/ERK: Cascata Canônica}
\begin{center}
\begin{tikzpicture}[node distance=1.5cm, scale=0.85, transform shape,
    box/.style={rectangle, draw, fill=azulclaro!20, text width=2.5cm, text centered, rounded corners, minimum height=0.8cm}]

    % Receptor
    \node[box, fill=laranjacel!30] (rtk) {RTK\\(EGFR, PDGFR)};

    % Adapter proteins
    \node[box, fill=roxodna!20, below of=rtk, node distance=1.2cm] (grb) {Grb2/SOS};

    % Ras
    \node[box, fill=verdeacido!30, below of=grb, node distance=1.2cm] (ras) {Ras-GTP};

    % MAPKKK
    \node[box, fill=azulclaro!40, below of=ras, node distance=1.2cm] (raf) {Raf\\(MAPKKK)};

    % MAPKK
    \node[box, fill=azulclaro!50, below of=raf, node distance=1.2cm] (mek) {MEK\\(MAPKK)};

    % MAPK
    \node[box, fill=azulclaro!60, below of=mek, node distance=1.2cm] (erk) {ERK\\(MAPK)};

    % Transcription factors
    \node[box, fill=yellow!30, below of=erk, node distance=1.2cm] (tf) {Fatores de\\Transcrição\\(c-Fos, Elk-1)};

    % Arrows
    \draw[->, thick] (rtk) -- (grb);
    \draw[->, thick] (grb) -- node[right, font=\tiny] {ativação} (ras);
    \draw[->, thick] (ras) -- (raf);
    \draw[->, thick] (raf) -- node[right, font=\tiny] {P} (mek);
    \draw[->, thick] (mek) -- node[right, font=\tiny] {P P} (erk);
    \draw[->, thick] (erk) -- (tf);

    % Response
    \node[below of=tf, node distance=1.2cm, text width=3cm, align=center, font=\small] (resp) {\textbf{Resposta:}\\Proliferação\\Diferenciação\\Sobrevivência};
\end{tikzpicture}
\end{center}
\end{frame}

\begin{frame}{Via PI3K/Akt: Sobrevivência e Metabolismo}
\begin{columns}
\column{0.5\textwidth}
\begin{center}
\begin{tikzpicture}[node distance=1.2cm, scale=0.75, transform shape,
    box/.style={rectangle, draw, text width=2cm, text centered, rounded corners, minimum height=0.7cm, font=\small}]

    \node[box, fill=laranjacel!30] (rtk) {RTK/GPCR};
    \node[box, fill=azulclaro!30, below of=rtk] (pi3k) {PI3K};
    \node[box, fill=verdeacido!30, below of=pi3k] (pip3) {PIP$_3$};
    \node[box, fill=roxodna!30, below of=pip3] (pdk) {PDK1};
    \node[box, fill=yellow!30, below of=pdk] (akt) {Akt (PKB)};

    \draw[->, thick] (rtk) -- (pi3k);
    \draw[->, thick] (pi3k) -- node[right, font=\tiny] {PIP$_2$ $\rightarrow$ PIP$_3$} (pip3);
    \draw[->, thick] (pip3) -- (pdk);
    \draw[->, thick] (pdk) -- node[right, font=\tiny] {P} (akt);

    % PTEN
    \node[box, fill=red!20, right of=pip3, node distance=3cm, text width=1.5cm] (pten) {PTEN\\(fosfatase)};
    \draw[-|, thick, red] (pten) -- (pip3);
\end{tikzpicture}
\end{center}

\column{0.5\textwidth}
\textbf{Alvos de Akt:}
\begin{itemize}
    \item \textbf{mTOR:} síntese proteica $\uparrow$
    \item \textbf{GSK3:} glicogênio síntese $\uparrow$
    \item \textbf{Bad:} apoptose $\downarrow$
    \item \textbf{FoxO:} gliconeogênese $\downarrow$
    \item \textbf{MDM2:} p53 $\downarrow$
\end{itemize}

\vspace{0.3cm}

\begin{exampleblock}{Relevância Clínica}
\begin{itemize}
    \item \textbf{PTEN:} supressor tumoral frequentemente mutado
    \item \textbf{PI3K/Akt:} hiperativa em $\sim$30\% dos cânceres
    \item \textbf{Diabetes:} resistência à insulina via Akt
\end{itemize}
\end{exampleblock}
\end{columns}
\end{frame}

\begin{frame}{Via JAK-STAT: Sinalização de Citocinas}
\begin{center}
\begin{tikzpicture}[node distance=1.5cm, scale=0.8, transform shape]
    % Membrane
    \draw[thick, double] (-2,0) -- (6,0);
    \node[left, font=\tiny] at (-2,0.3) {Extracelular};
    \node[left, font=\tiny] at (-2,-0.3) {Intracelular};

    % Cytokine receptor dimer
    \draw[thick, azulclaro, fill=azulclaro!30] (0,-0.3) rectangle (0.5,1);
    \draw[thick, azulclaro, fill=azulclaro!30] (1.5,-0.3) rectangle (2,1);

    % Cytokine
    \draw[fill=laranjacel!50] (1,1.5) circle (0.3);
    \node[above, font=\small] at (1,1.8) {Citocina};

    % JAK proteins
    \draw[thick, verdeacido, fill=verdeacido!30] (0,-0.8) circle (0.4);
    \draw[thick, verdeacido, fill=verdeacido!30] (2,-0.8) circle (0.4);
    \node[font=\small] at (0,-0.8) {JAK};
    \node[font=\small] at (2,-0.8) {JAK};

    % Phosphorylations
    \foreach \x in {0, 2} {
        \foreach \y in {-1.4, -1.6, -1.8} {
            \node[font=\scriptsize, red] at (\x,\y) {P};
        }
    }

    % STAT proteins
    \draw[thick, roxodna, fill=roxodna!30] (1,-2.3) ellipse (0.5 and 0.3);
    \node[font=\small] at (1,-2.3) {STAT};
    \draw[->, thick] (1,-1.9) -- (1,-2);

    % Nucleus
    \draw[thick, dashed] (3.5,-3) rectangle (5.5,-4.5);
    \node[font=\small] at (4.5,-3.2) {Núcleo};

    % STAT dimer in nucleus
    \draw[thick, yellow, fill=yellow!30] (4.5,-3.9) ellipse (0.6 and 0.3);
    \node[font=\small] at (4.5,-3.9) {STAT};
    \node[font=\tiny, red] at (4.5,-3.6) {P P};

    % Arrow to nucleus
    \draw[->, thick, bend right=30] (1.5,-2.5) to node[below, font=\tiny] {dimerização} (3.5,-3.8);

    % DNA
    \draw[thick, azulescuro] (4,-4.3) -- (5,-4.3);
    \node[below, font=\tiny] at (4.5,-4.5) {Gene alvo};
\end{tikzpicture}
\end{center}

\importante{
Via rápida e direta: do receptor à transcrição gênica sem cascatas intermediárias complexas!
}
\end{frame}

\begin{frame}{Via Wnt: Desenvolvimento e Polaridade}
\begin{columns}
\column{0.5\textwidth}
\textbf{Via Canônica Wnt/$\beta$-catenina:}

\begin{enumerate}
    \item \textbf{Sem Wnt:}
    \begin{itemize}
        \item $\beta$-catenina fosforilada por complexo de destruição (GSK3, APC, Axina)
        \item Ubiquitinação e degradação
    \end{itemize}

    \vspace{0.2cm}

    \item \textbf{Com Wnt:}
    \begin{itemize}
        \item Wnt liga a Frizzled + LRP5/6
        \item Dishevelled (Dvl) ativado
        \item Complexo de destruição inibido
        \item $\beta$-catenina acumula
        \item Entra no núcleo
        \item Ativa transcrição (TCF/LEF)
    \end{itemize}
\end{enumerate}

\column{0.5\textwidth}
\begin{center}
\begin{tikzpicture}[scale=0.6]
    % OFF state
    \begin{scope}
        \node[font=\small, align=center] at (0,3) {\textbf{OFF}\\(sem Wnt)};

        % Destruction complex
        \draw[thick, fill=red!20] (0,1.5) ellipse (1 and 0.6);
        \node[font=\tiny, align=center] at (0,1.5) {Complexo\\destruição\\GSK3/APC};

        % beta-catenin degraded
        \draw[thick, fill=azulclaro!30] (0,0.5) rectangle (0.8,0.9);
        \node[font=\tiny] at (0.4,0.7) {$\beta$-cat};
        \node[red, font=\tiny] at (0.4,0.3) {P};
        \draw[->, thick, red] (0.4,0.2) -- (0.4,-0.3);
        \node[font=\tiny] at (0.4,-0.6) {Degradação};
    \end{scope}

    % ON state
    \begin{scope}[xshift=4cm]
        \node[font=\small, align=center] at (0,3) {\textbf{ON}\\(com Wnt)};

        % Wnt ligand
        \draw[fill=laranjacel!50] (0,2.3) circle (0.2);
        \node[font=\tiny, above] at (0,2.5) {Wnt};

        % Receptors
        \draw[thick, azulclaro] (0,1.5) rectangle (0.3,2.2);

        % Destruction complex inhibited
        \draw[thick, fill=gray!20] (0,0.8) ellipse (1 and 0.6);
        \node[font=\tiny, align=center] at (0,0.8) {Complexo\\inativo};

        % beta-catenin accumulates
        \draw[thick, fill=verdeacido!30] (-0.5,0) rectangle (-0.1,0.4);
        \draw[thick, fill=verdeacido!30] (0.1,0) rectangle (0.5,0.4);
        \node[font=\tiny] at (0,0.2) {$\beta$-cat};

        % Nucleus
        \draw[thick, dashed] (-0.8,-1) rectangle (0.8,-2.2);
        \node[font=\tiny] at (0,-1.2) {Núcleo};

        % beta-catenin in nucleus
        \draw[thick, fill=yellow!30] (0,-1.7) rectangle (0.4,-1.5);
        \node[font=\tiny] at (0.2,-1.6) {$\beta$-cat};

        \draw[->, thick] (0.2,-0.2) -- (0.2,-1);
    \end{scope}
\end{tikzpicture}
\end{center}
\end{columns}

\vspace{0.2cm}

\exemplo{Wnt em Câncer}{
Mutações em APC (via Wnt) são encontradas em $>$80\% dos cânceres colorretais.
}
\end{frame}

\begin{frame}{Via Notch: Decisões de Destino Celular}
\begin{center}
\begin{tikzpicture}[scale=0.8]
    % Cell 1 (sending)
    \begin{scope}
        \draw[thick] (-2,-1) rectangle (0,3);
        \node[font=\small] at (-1,2.5) {Célula Sinalizadora};

        % Delta ligand
        \draw[thick, fill=laranjacel!40] (-0.3,1) -- (0.3,1) -- (0.3,1.8) -- (-0.3,1.8) -- cycle;
        \node[font=\tiny, right] at (0.3,1.4) {Delta};
    \end{scope}

    % Cell 2 (receiving)
    \begin{scope}[xshift=4cm]
        \draw[thick] (-2,-1) rectangle (0,3);
        \node[font=\small] at (-1,2.5) {Célula Receptora};

        % Notch receptor
        \draw[thick, fill=azulclaro!40] (-2.3,0.8) -- (-2.3,2) -- (-1.5,2) -- (-1.5,0.8);
        \node[font=\tiny, left] at (-2.3,1.4) {Notch};

        % Extracellular domain (NECD)
        \draw[thick, fill=roxodna!30] (-2.7,1.5) circle (0.3);
        \node[font=\tiny] at (-2.7,1.5) {NECD};

        % Intracellular domain (NICD)
        \draw[thick, fill=verdeacido!30] (-1.3,1.2) ellipse (0.4 and 0.3);
        \node[font=\tiny] at (-1.3,1.2) {NICD};

        % Cleavage
        \draw[->, thick, red] (-2.7,1.1) -- (-3.2,0.5);
        \node[font=\tiny, red] at (-3.2,0.3) {clivagem};

        % NICD to nucleus
        \draw[->, thick] (-1,1) -- (-1,0);

        % Nucleus
        \draw[thick, dashed] (-1.8,-0.3) rectangle (-0.2,-1.8);
        \node[font=\tiny] at (-1,-0.5) {Núcleo};

        % NICD + CSL
        \draw[thick, fill=yellow!30] (-1.5,-1) rectangle (-0.5,-1.3);
        \node[font=\tiny] at (-1,-1.15) {NICD+CSL};

        % DNA
        \draw[thick, azulescuro] (-1.7,-1.5) -- (-0.3,-1.5);
        \node[font=\tiny, below] at (-1,-1.7) {Gene alvo};
    \end{scope}

    % Interaction
    \draw[<->, thick, laranjacel] (0.3,1.4) -- (1.7,1.4);
    \node[font=\tiny, above] at (1,1.4) {Interação};
\end{tikzpicture}
\end{center}

\importante{
Notch é único: o receptor é clivado e o fragmento intracelular vai diretamente ao núcleo!
}
\end{frame}

\begin{frame}{Via TGF-$\beta$/Smad}
\begin{columns}
\column{0.5\textwidth}
\begin{center}
\begin{tikzpicture}[node distance=1.2cm, scale=0.7, transform shape,
    box/.style={rectangle, draw, text width=2.2cm, text centered, rounded corners, minimum height=0.7cm, font=\small}]

    % Membrane
    \draw[thick, double] (-1,0) -- (3,0);

    % TGF-beta ligand
    \draw[fill=laranjacel!50] (1,0.8) circle (0.3);
    \node[above, font=\tiny] at (1,1.1) {TGF-$\beta$};

    % Type II and Type I receptors
    \draw[thick, fill=azulclaro!30] (0.3,-0.2) rectangle (0.8,0.5);
    \draw[thick, fill=azulclaro!30] (1.2,-0.2) rectangle (1.7,0.5);
    \node[font=\tiny] at (0.55,-0.5) {T$\beta$RII};
    \node[font=\tiny] at (1.45,-0.5) {T$\beta$RI};

    % R-Smads
    \node[box, fill=verdeacido!30] at (1,-1.5) {R-Smad\\(Smad2/3)};

    % Co-Smad
    \node[box, fill=roxodna!30] at (1,-2.5) {Co-Smad\\(Smad4)};

    % Complex
    \node[box, fill=yellow!30] at (1,-3.5) {Complexo\\Smad};

    % Nucleus
    \draw[thick, dashed] (0,-4.5) rectangle (2,-5.5);
    \node[font=\tiny] at (1,-4.7) {Núcleo};

    % Transcription
    \node[font=\tiny, align=center] at (1,-5.2) {Transcrição\\Gene alvo};

    % Arrows
    \draw[->, thick] (1,-0.7) -- (1,-1.2);
    \draw[->, thick] (1,-1.8) -- (1,-2.2);
    \draw[->, thick] (1,-2.8) -- (1,-3.2);
    \draw[->, thick] (1,-3.8) -- (1,-4.3);
\end{tikzpicture}
\end{center}

\column{0.5\textwidth}
\textbf{Regulação:}
\begin{itemize}
    \item \textbf{I-Smads} (Smad6/7): inibem R-Smads (feedback negativo)
    \item \textbf{Ubiquitinação:} degradação de R-Smads
    \item \textbf{Fosfatases:} desfosforilam Smads
\end{itemize}

\vspace{0.3cm}

\textbf{Funções Biológicas:}
\begin{itemize}
    \item Diferenciação celular
    \item Apoptose
    \item Transição epitélio-mesênquima (EMT)
    \item Fibrose
    \item Imunossupressão
\end{itemize}

\vspace{0.3cm}

\begin{exampleblock}{Dual Role}
TGF-$\beta$ age como \textbf{supressor tumoral} em estágios iniciais, mas promove \textbf{metástase} em estágios avançados.
\end{exampleblock}
\end{columns}
\end{frame}

\begin{frame}{Crosstalk entre Vias de Sinalização}
\begin{center}
\begin{tikzpicture}[scale=0.85, every node/.style={font=\small}]
    % Central pathways as nodes
    \node[draw, circle, fill=azulclaro!30, text width=1.5cm, align=center] (mapk) at (0,0) {MAPK/\\ERK};
    \node[draw, circle, fill=verdeacido!30, text width=1.5cm, align=center] (pi3k) at (3,1.5) {PI3K/\\Akt};
    \node[draw, circle, fill=roxodna!30, text width=1.5cm, align=center] (jak) at (3,-1.5) {JAK-\\STAT};
    \node[draw, circle, fill=laranjacel!30, text width=1.5cm, align=center] (wnt) at (6,0) {Wnt};
    \node[draw, circle, fill=yellow!30, text width=1.5cm, align=center] (tgf) at (4.5,-3) {TGF-$\beta$};
    \node[draw, circle, fill=red!20, text width=1.5cm, align=center] (nfkb) at (1.5,-3) {NF-$\kappa$B};

    % Crosstalk arrows
    \draw[<->, thick, azulescuro] (mapk) -- node[above, font=\tiny] {Ras-PI3K} (pi3k);
    \draw[<->, thick, azulescuro] (mapk) -- node[below, font=\tiny] {ERK-STAT} (jak);
    \draw[<->, thick, azulescuro] (pi3k) -- node[right, font=\tiny] {GSK3} (wnt);
    \draw[<->, thick, azulescuro] (pi3k) -- (jak);
    \draw[<->, thick, azulescuro] (jak) -- (nfkb);
    \draw[<->, thick, azulescuro] (mapk) -- (nfkb);
    \draw[<->, thick, azulescuro] (tgf) -- (mapk);
    \draw[<->, thick, azulescuro] (tgf) -- (wnt);
\end{tikzpicture}
\end{center}

\importante{
O crosstalk permite:
\begin{itemize}
    \item Integração de múltiplos sinais
    \item Ajuste fino da resposta celular
    \item Comportamentos complexos (biestabilidade, oscilações)
    \item Robustez e redundância
\end{itemize}
}
\end{frame}

% ============================================================================
% SEÇÃO 4: CASCATAS DE QUINASES
% ============================================================================
\section{Cascatas de Proteínas Quinases}

\begin{frame}{9.4 Cascatas de Fosforilação}
\definicao{Cascata de Quinases}{
Sequência de reações onde uma quinase ativa outra por fosforilação, resultando em amplificação de sinal.
}

\begin{center}
\begin{tikzpicture}[node distance=1.5cm, scale=0.9, transform shape]
    % Generic cascade
    \node[draw, rectangle, fill=azulclaro!30, text width=2.5cm, align=center] (k1) {Quinase 1\\(inativa)};
    \node[draw, rectangle, fill=azulclaro!50, text width=2.5cm, align=center, right of=k1, node distance=4cm] (k1a) {Quinase 1*\\(ativa)};

    \node[draw, rectangle, fill=verdeacido!30, text width=2.5cm, align=center, below of=k1a] (k2) {Quinase 2\\(inativa)};
    \node[draw, rectangle, fill=verdeacido!50, text width=2.5cm, align=center, right of=k2, node distance=4cm] (k2a) {Quinase 2**\\(ativa)};

    \node[draw, rectangle, fill=laranjacel!30, text width=2.5cm, align=center, below of=k2a] (k3) {Quinase 3\\(inativa)};
    \node[draw, rectangle, fill=laranjacel!50, text width=2.5cm, align=center, right of=k3, node distance=4cm] (k3a) {Quinase 3***\\(ativa)};

    % Arrows
    \draw[->, thick, red] (k1) -- node[above, font=\small] {sinal} (k1a);
    \draw[->, thick] (k1a) -- node[right, font=\small] {P} (k2);
    \draw[->, thick] (k2) -- (k2a);
    \draw[->, thick] (k2a) -- node[right, font=\small] {P P} (k3);
    \draw[->, thick] (k3) -- (k3a);

    % Amplification labels
    \node[right, font=\small, text width=2cm] at (9.5,0) {1 molécula};
    \node[right, font=\small, text width=2cm] at (9.5,-1.5) {10-100\\moléculas};
    \node[right, font=\small, text width=2cm] at (9.5,-3) {100-1000\\moléculas};
\end{tikzpicture}
\end{center}

\importante{
Cascatas proporcionam \textbf{amplificação}, \textbf{ultrassensibilidade} e \textbf{integração} de sinais!
}
\end{frame}

\begin{frame}{Cascata MAPK: MAPKKK $\rightarrow$ MAPKK $\rightarrow$ MAPK}
\begin{block}{Estrutura da Cascata}
\begin{equation*}
\text{MAPKKK} \xrightarrow{P} \text{MAPKK} \xrightarrow{P \, P} \text{MAPK} \xrightarrow{P} \text{Substrato}
\end{equation*}

\begin{itemize}
    \item \textbf{MAPKKK:} Raf, MEKK, MLK, TAK1
    \item \textbf{MAPKK:} MEK1/2, MKK3/6, MKK4/7
    \item \textbf{MAPK:} ERK1/2, p38, JNK
\end{itemize}
\end{block}

\begin{exampleblock}{Dupla Fosforilação de MAPK}
MAPK requer fosforilação em \textbf{dois} resíduos (Thr e Tyr no motif TXY) para ativação completa:
\begin{equation*}
\text{MAPK} \xrightarrow{\text{MAPKK}} \text{MAPK-P} \xrightarrow{\text{MAPKK}} \text{MAPK-PP*}
\end{equation*}

Isso confere \textbf{ultrassensibilidade}!
\end{exampleblock}
\end{frame}

\begin{frame}{Ultrassensibilidade: Função de Goldbeter-Koshland}
\begin{columns}
\column{0.5\textwidth}
\begin{block}{Modelo de GK}
Para um ciclo de modificação (fosforilação/desfosforilação):

\begin{equation*}
\frac{[X^*]}{[X]_{total}} = \frac{2vq}{v + q - v + q - (v - q)^2 + 4vq}
\end{equation*}

Onde:
\begin{itemize}
    \item $v = \frac{V_1[S]}{K_1 + [S]}$ (quinase)
    \item $q = \frac{V_2}{K_2}$ (fosfatase)
\end{itemize}

Quando $K_1, K_2 \ll [X]_{total}$: comportamento \textbf{ultrassensível}
\end{block}

\column{0.5\textwidth}
\textbf{Coeficiente de Hill Efetivo:}
\begin{equation*}
n_{H} \approx 1 + \frac{2[X]_{total}}{K_1 + K_2}
\end{equation*}

\vspace{0.3cm}

\begin{exampleblock}{Interpretação}
\begin{itemize}
    \item $n_H = 1$: resposta hiperbólica (Michaelis-Menten)
    \item $n_H > 1$: ultrassensível (switch-like)
    \item Cascatas de MAPK: $n_H$ pode chegar a 5-10!
\end{itemize}
\end{exampleblock}
\end{columns}

\vspace{0.3cm}

\importante{
Ultrassensibilidade permite decisões binárias (ON/OFF) mesmo com sinais graduais!
}
\end{frame}

\begin{frame}{Amplificação de Sinal em Cascatas}
\begin{block}{Modelo de Amplificação}
Para uma cascata de $n$ níveis, cada nível com ganho $g_i$:

\begin{equation*}
\text{Amplificação Total} = \prod_{i=1}^{n} g_i = g_1 \times g_2 \times \cdots \times g_n
\end{equation*}

\textbf{Ganho em cada nível:}
\begin{equation*}
g_i = \frac{k_{cat,i}}{k_{inact,i}} \times \frac{[K_i]_{total}}{[K_{i-1}]_{ativo}}
\end{equation*}
\end{block}

\exemplo{Cascata MAPK Típica}{
\begin{itemize}
    \item Raf: 10 moléculas ativas
    \item MEK: 100 moléculas ativas (ganho = 10)
    \item ERK: 1000 moléculas ativas (ganho = 10)
    \item \textbf{Amplificação total:} $10 \times 10 = 100\times$
\end{itemize}
}
\end{frame}

\begin{frame}{Fosforilação Distributiva vs Processiva}
\begin{columns}
\column{0.5\textwidth}
\begin{block}{Distributiva}
Substrato se dissocia após cada fosforilação:

\vspace{0.2cm}

\begin{tikzpicture}[scale=0.6]
    \node[draw, fill=azulclaro!30] (s) at (0,0) {S};
    \node[draw, fill=verdeacido!30] (k) at (2,0) {K};
    \node[draw, fill=azulclaro!50] (sp) at (0,-1.5) {SP};
    \node[draw, fill=azulclaro!70] (spp) at (0,-3) {SPP};

    \draw[->, thick] (s) -- (k);
    \draw[->, thick] (k) -- (sp);
    \draw[->, thick, dashed, red] (sp) -- (0,-1) node[right, font=\tiny] {dissocia};
    \draw[->, thick] (sp) -- (k);
    \draw[->, thick] (k) -- (spp);
\end{tikzpicture}

\begin{itemize}
    \item Mais comum
    \item Gera \textbf{ultrassensibilidade}
    \item Coeficiente de Hill alto
\end{itemize}
\end{block}

\column{0.5\textwidth}
\begin{block}{Processiva}
Substrato permanece ligado:

\vspace{0.2cm}

\begin{tikzpicture}[scale=0.6]
    \node[draw, fill=azulclaro!30] (s) at (0,0) {S};
    \node[draw, fill=verdeacido!30] (k) at (2,0) {K};
    \node[draw, fill=azulclaro!50] (sp) at (2,-1.5) {S-K-P};
    \node[draw, fill=azulclaro!70] (spp) at (2,-3) {SPP};

    \draw[->, thick] (s) -- (k);
    \draw[->, thick] (k) -- node[right, font=\tiny] {P} (sp);
    \draw[->, thick] (sp) -- node[right, font=\tiny] {P} (spp);
    \draw[->, thick] (spp) -- (0,-3);
\end{tikzpicture}

\begin{itemize}
    \item Menos comum
    \item Resposta mais gradual
    \item Coeficiente de Hill $\approx 1$
\end{itemize}
\end{block}
\end{columns}

\vspace{0.3cm}

\begin{exampleblock}{MAPK é Distributiva}
MEK fosforila ERK de forma distributiva $\rightarrow$ ultrassensibilidade!
\end{exampleblock}
\end{frame}

\begin{frame}{Biestabilidade em Cascatas de Sinalização}
\begin{columns}
\column{0.5\textwidth}
\definicao{Biestabilidade}{
Sistema com dois estados estáveis (ON e OFF) separados por ponto de instabilidade.
}

\textbf{Requisitos:}
\begin{itemize}
    \item Ultrassensibilidade ($n_H > 1$)
    \item Feedback positivo
    \item Saturação de fosfatases
\end{itemize}

\vspace{0.3cm}

\textbf{Consequências:}
\begin{itemize}
    \item Memória celular
    \item Decisões irreversíveis
    \item Histerese
    \item Resistência a ruído
\end{itemize}

\column{0.5\textwidth}
\begin{center}
\begin{tikzpicture}[scale=0.8]
    % Axes
    \draw[->] (0,0) -- (5,0) node[right, font=\small] {Estímulo};
    \draw[->] (0,0) -- (0,4) node[above, font=\small] {Resposta};

    % Bistable curve (hysteresis)
    \draw[thick, azulescuro] (0.5,0.5) -- (2,0.7);
    \draw[thick, azulescuro] (2,0.7) -- (2.2,2.5);
    \draw[thick, azulescuro] (2.2,2.5) -- (4.5,3.5);

    \draw[thick, red, dashed] (0.5,0.5) -- (3,1);
    \draw[thick, red, dashed] (3,1) -- (3.2,2.8);
    \draw[thick, red, dashed] (3.2,2.8) -- (4.5,3.5);

    % Labels
    \node[font=\small] at (1,0.3) {OFF};
    \node[font=\small] at (4,3.8) {ON};
    \node[font=\tiny, red] at (2.6,1.8) {Instável};

    % Arrows showing hysteresis
    \draw[->, thick, verdeacido] (1.5,1.5) to [out=30, in=150] (3.5,2);
    \draw[->, thick, laranjacel] (3.5,2.5) to [out=210, in=330] (1.5,2);

    \node[font=\tiny] at (2.5,0.5) {$S_{ON}$};
    \node[font=\tiny] at (3.5,0.5) {$S_{OFF}$};
\end{tikzpicture}
\end{center}

\begin{exampleblock}{Histerese}
Threshold para ativar $\neq$ threshold para desativar!
\end{exampleblock}
\end{columns}
\end{frame}

\begin{frame}[fragile]{Simulação de Cascata MAPK}
\begin{lstlisting}[language=Python, caption=Modelo simplificado de cascata MAPK, basicstyle=\ttfamily\scriptsize]
import numpy as np
from scipy.integrate import odeint
import matplotlib.pyplot as plt

def mapk_cascade(y, t, stimulus):
    """Modelo de cascata MAPK de 3 niveis"""
    K1, K1p, K2, K2p, K3, K3p = y

    # Parametros
    k1 = 0.1 * stimulus  # ativacao de K1
    k2 = 0.5              # K1p ativa K2
    k3 = 1.0              # K2p ativa K3
    d1 = 0.1              # desfosforilacao
    d2 = 0.2
    d3 = 0.3

    # ODEs
    dK1 = -k1*K1 + d1*K1p
    dK1p = k1*K1 - d1*K1p

    dK2 = -k2*K1p*K2 + d2*K2p
    dK2p = k2*K1p*K2 - d2*K2p

    dK3 = -k3*K2p*K3 + d3*K3p
    dK3p = k3*K2p*K3 - d3*K3p

    return [dK1, dK1p, dK2, dK2p, dK3, dK3p]

# Condicoes iniciais (total = 100 para cada nivel)
y0 = [100, 0, 100, 0, 100, 0]
t = np.linspace(0, 50, 500)
stimulus = 1.0

sol = odeint(mapk_cascade, y0, t, args=(stimulus,))
\end{lstlisting}
\end{frame}

\begin{frame}[fragile]{Simulação de Cascata MAPK (cont.)}
\begin{lstlisting}[language=Python, caption=Visualização da cascata, basicstyle=\ttfamily\scriptsize]
# Plotar resultados
plt.figure(figsize=(12, 6))

plt.subplot(1, 2, 1)
plt.plot(t, sol[:, 1], 'b-', linewidth=2, label='K1* (MAPKKK)')
plt.plot(t, sol[:, 3], 'g-', linewidth=2, label='K2* (MAPKK)')
plt.plot(t, sol[:, 5], 'r-', linewidth=2, label='K3* (MAPK)')
plt.xlabel('Tempo')
plt.ylabel('Concentracao (forma ativa)')
plt.title('Dinamica Temporal da Cascata MAPK')
plt.legend()
plt.grid(True, alpha=0.3)

# Dose-response
plt.subplot(1, 2, 2)
stimuli = np.logspace(-2, 1, 20)
responses = []

for stim in stimuli:
    sol = odeint(mapk_cascade, y0, t, args=(stim,))
    responses.append(sol[-1, 5])  # K3p no equilibrio

plt.semilogx(stimuli, responses, 'ro-', linewidth=2)
plt.xlabel('Estimulo (log scale)')
plt.ylabel('MAPK ativa (equilibrio)')
plt.title('Curva Dose-Resposta (Ultrassensibilidade)')
plt.grid(True, alpha=0.3)
plt.tight_layout()
\end{lstlisting}
\end{frame}

% ============================================================================
% SEÇÃO 5: SEGUNDOS MENSAGEIROS
% ============================================================================
\section{Segundos Mensageiros}

\begin{frame}{9.5 Segundos Mensageiros}
\definicao{Segundo Mensageiro}{
Molécula intracelular pequena que propaga o sinal de um receptor para efetores intracelulares, amplificando e difundindo o sinal.
}

\begin{block}{Principais Segundos Mensageiros}
\begin{itemize}
    \item \textbf{cAMP} (AMP cíclico)
    \item \textbf{cGMP} (GMP cíclico)
    \item \textbf{Ca}$^{2+}$ (cálcio)
    \item \textbf{IP}$_3$ (inositol 1,4,5-trifosfato)
    \item \textbf{DAG} (diacilglicerol)
    \item \textbf{NO} (óxido nítrico)
\end{itemize}
\end{block}

\importante{
Segundos mensageiros permitem:
\begin{itemize}
    \item \textbf{Amplificação:} 1 receptor $\rightarrow$ milhares de moléculas de 2º mensageiro
    \item \textbf{Difusão:} sinal se espalha pelo citoplasma
    \item \textbf{Integração:} múltiplos sinais convergem
\end{itemize}
}
\end{frame}

\begin{frame}{cAMP: Sinalização via Adenilil Ciclase}
\begin{columns}
\column{0.5\textwidth}
\textbf{Via de Síntese:}
\begin{enumerate}
    \item GPCR ativa G$_s$-$\alpha$
    \item G$_s$-$\alpha$-GTP ativa adenilil ciclase (AC)
    \item AC converte ATP $\rightarrow$ cAMP
    \item cAMP ativa PKA (proteína quinase A)
    \item PKA fosforila alvos downstream
    \item Fosfodiesterase (PDE) degrada cAMP $\rightarrow$ AMP
\end{enumerate}

\vspace{0.3cm}

\textbf{Alvos de PKA:}
\begin{itemize}
    \item CREB (transcrição)
    \item Enzimas metabólicas
    \item Canais iônicos
\end{itemize}

\column{0.5\textwidth}
\begin{center}
\begin{tikzpicture}[node distance=1.2cm, scale=0.7, transform shape,
    box/.style={rectangle, draw, text width=2cm, text centered, rounded corners, minimum height=0.7cm, font=\small}]

    \node[box, fill=laranjacel!30] (gpcr) {GPCR};
    \node[box, fill=verdeacido!30, below of=gpcr] (gs) {G$_s$-$\alpha$};
    \node[box, fill=azulclaro!30, below of=gs] (ac) {Adenilil\\Ciclase};
    \node[box, fill=yellow!30, below of=ac] (camp) {cAMP};
    \node[box, fill=roxodna!30, below of=camp] (pka) {PKA\\(ativa)};
    \node[box, fill=red!20, below of=pka] (target) {Alvos\\(fosforilados)};

    \draw[->, thick] (gpcr) -- (gs);
    \draw[->, thick] (gs) -- (ac);
    \draw[->, thick] (ac) -- node[right, font=\tiny] {ATP} (camp);
    \draw[->, thick] (camp) -- (pka);
    \draw[->, thick] (pka) -- node[right, font=\tiny] {P} (target);

    % PDE
    \draw[-|, thick, red] (3,-3) node[right, font=\tiny] {PDE} -- (camp);
\end{tikzpicture}
\end{center}

\begin{exampleblock}{Exemplos}
\begin{itemize}
    \item Epinefrina $\rightarrow$ glicogenólise
    \item ACTH $\rightarrow$ síntese de cortisol
    \item Glucagon $\rightarrow$ liberação de glicose
\end{itemize}
\end{exampleblock}
\end{columns}
\end{frame}

\begin{frame}{Sinalização por Ca$^{2+}$: Oscilações e Ondas}
\begin{columns}
\column{0.5\textwidth}
\textbf{Fontes de Ca$^{2+}$:}
\begin{itemize}
    \item \textbf{Intracelular:} retículo endoplasmático (ER)
    \item \textbf{Extracelular:} através de canais
\end{itemize}

\vspace{0.3cm}

\textbf{Concentrações:}
\begin{itemize}
    \item [Ca$^{2+}$]$_{cito,repouso}$ $\sim$ 100 nM
    \item [Ca$^{2+}$]$_{ER}$ $\sim$ 0.5 mM
    \item [Ca$^{2+}$]$_{extra}$ $\sim$ 1-2 mM
\end{itemize}

\vspace{0.3cm}

\textbf{Sensores de Ca$^{2+}$:}
\begin{itemize}
    \item \textbf{Calmodulina (CaM)}
    \item \textbf{Troponina C} (músculo)
    \item \textbf{PKC} (proteína quinase C)
    \item \textbf{CaMKII}
\end{itemize}

\column{0.5\textwidth}
\begin{center}
\begin{tikzpicture}[scale=0.7]
    % Axes
    \draw[->] (0,0) -- (6,0) node[right, font=\small] {Tempo};
    \draw[->] (0,0) -- (0,4) node[above, font=\small] {[Ca$^{2+}$]};

    % Baseline
    \draw[thick, gray, dashed] (0,0.5) -- (6,0.5);
    \node[left, font=\tiny] at (0,0.5) {basal};

    % Oscillations
    \draw[thick, azulescuro, domain=0:6, samples=100, smooth]
        plot (\x, {0.5 + 2.5*exp(-0.2*\x)*sin(360*\x)});

    \node[font=\small, align=center] at (3,-1) {Oscilações de Ca$^{2+}$};
\end{tikzpicture}
\end{center}

\vspace{0.5cm}

\begin{exampleblock}{Padrões de Ca$^{2+}$}
\begin{itemize}
    \item \textbf{Transiente:} pico único
    \item \textbf{Oscilações:} picos periódicos
    \item \textbf{Ondas:} propagação espacial
\end{itemize}

Frequência e amplitude codificam informação!
\end{exampleblock}
\end{columns}
\end{frame}

\begin{frame}{Via IP$_3$/DAG: Fosfolipase C}
\begin{center}
\begin{tikzpicture}[node distance=1.5cm, scale=0.85, transform shape]
    % Membrane
    \draw[thick, double] (-2,0) -- (6,0);
    \node[left, font=\tiny] at (-2,0.3) {Extracelular};
    \node[left, font=\tiny] at (-2,-0.3) {Intracelular};

    % GPCR
    \draw[thick, fill=laranjacel!30] (0,-0.3) rectangle (0.5,0.5);
    \node[above, font=\tiny] at (0.25,0.7) {GPCR};

    % Gq-alpha
    \draw[thick, fill=verdeacido!30] (1,-0.7) circle (0.4);
    \node[font=\small] at (1,-0.7) {G$_q$};

    % PLC
    \draw[thick, fill=azulclaro!30] (2.5,-0.7) ellipse (0.6 and 0.4);
    \node[font=\small] at (2.5,-0.7) {PLC};

    % PIP2 in membrane
    \draw[thick, fill=roxodna!30] (2.5,0) circle (0.2);
    \node[above, font=\tiny] at (2.5,0.3) {PIP$_2$};

    % IP3 and DAG
    \node[draw, rectangle, fill=yellow!30, rounded corners] (ip3) at (4.5,-1.5) {IP$_3$};
    \node[draw, rectangle, fill=red!20, rounded corners] (dag) at (2.5,-1.5) {DAG};

    % Arrows
    \draw[->, thick] (0.25,-0.3) -- (1,-0.3);
    \draw[->, thick] (1.4,-0.7) -- (1.9,-0.7);
    \draw[->, thick] (2.5,-0.3) -- (2.5,-1);
    \draw[->, thick] (2.7,-1.1) -- (ip3);
    \draw[->, thick] (2.3,-1.1) -- (dag);

    % ER and Ca2+
    \draw[thick, dashed] (4,-3) rectangle (6,-4);
    \node[font=\tiny] at (5,-3.2) {ER};
    \node[font=\small] at (5,-3.6) {Ca$^{2+}$};

    % IP3 receptor
    \draw[->, thick, red] (ip3) -- (5,-2.8);
    \node[right, font=\tiny] at (5,-2.5) {IP$_3$R};

    % Ca2+ release
    \foreach \i in {1,2,3} {
        \draw[->, thick, red] (5,-2.8) -- ++(0.2*\i, 0.5*\i);
    }

    % PKC activation
    \node[draw, rectangle, fill=verdeacido!40, rounded corners] at (1,-3) {PKC};
    \draw[->, thick] (dag) -- (1,-2);
    \draw[->, thick] (4,-2) to [out=180, in=90] (1,-2.7);
    \node[left, font=\tiny] at (2.5,-2.5) {Ca$^{2+}$};
\end{tikzpicture}
\end{center}

\importante{
IP$_3$ e DAG atuam sinergisticamente: IP$_3$ libera Ca$^{2+}$, Ca$^{2+}$ + DAG ativam PKC!
}
\end{frame}

\begin{frame}{Modelo Matemático de Oscilações de Ca$^{2+}$}
\begin{block}{Modelo Simplificado}
\begin{equation*}
\begin{aligned}
\frac{d[Ca^{2+}]}{dt} &= v_{in} - v_{out} + v_{release} - v_{uptake}\\
v_{release} &= k_{rel} \cdot [IP_3] \cdot \frac{[Ca^{2+}]^2}{K_1^2 + [Ca^{2+}]^2}\\
v_{uptake} &= k_{uptake} \cdot [Ca^{2+}]
\end{aligned}
\end{equation*}

\textbf{CICR (Calcium-Induced Calcium Release):}
\begin{itemize}
    \item Ca$^{2+}$ estimula sua própria liberação (feedback positivo)
    \item Combinado com uptake (feedback negativo)
    \item Resultado: oscilações!
\end{itemize}
\end{block}

\begin{exampleblock}{Codificação por Frequência (Frequency Encoding)}
A \textbf{frequência} das oscilações de Ca$^{2+}$ codifica a intensidade do estímulo:
\begin{itemize}
    \item Estímulo fraco $\rightarrow$ baixa frequência
    \item Estímulo forte $\rightarrow$ alta frequência
\end{itemize}
\end{exampleblock}
\end{frame}

\begin{frame}[fragile]{Simulação de Oscilações de Ca$^{2+}$}
\begin{lstlisting}[language=Python, caption=Modelo de oscilações de cálcio, basicstyle=\ttfamily\scriptsize]
import numpy as np
from scipy.integrate import odeint
import matplotlib.pyplot as plt

def ca_oscillations(y, t, IP3):
    """Modelo simplificado de oscilacoes de Ca2+"""
    Ca = y[0]

    # Parametros
    k_rel = 2.0       # liberacao do ER
    k_uptake = 0.5    # recaptura para o ER
    K1 = 0.3          # constante para CICR
    v_in = 0.1        # influxo externo
    v_out = 0.05      # efluxo

    # Liberacao induzida por IP3 e Ca (CICR)
    v_release = k_rel * IP3 * (Ca**2 / (K1**2 + Ca**2))

    # Uptake pelo ER (SERCA pumps)
    v_uptake_rate = k_uptake * Ca

    # ODE
    dCa = v_in - v_out + v_release - v_uptake_rate

    return [dCa]

# Simular para diferentes niveis de IP3
t = np.linspace(0, 100, 1000)
IP3_levels = [0.5, 1.0, 1.5]

plt.figure(figsize=(12, 4))
for i, IP3 in enumerate(IP3_levels):
    Ca0 = [0.1]
    sol = odeint(ca_oscillations, Ca0, t, args=(IP3,))
    plt.subplot(1, 3, i+1)
    plt.plot(t, sol[:, 0], 'b-', linewidth=2)
    plt.xlabel('Tempo')
    plt.ylabel('[Ca2+]')
    plt.title(f'IP3 = {IP3}')
    plt.grid(True, alpha=0.3)
\end{lstlisting}
\end{frame}

\begin{frame}{cGMP e Sinalização por NO}
\begin{columns}
\column{0.5\textwidth}
\textbf{Via NO/cGMP:}
\begin{enumerate}
    \item NO (óxido nítrico) produzido por NOS (NO sintase)
    \item NO difunde através de membranas
    \item NO ativa guanilil ciclase solúvel (sGC)
    \item sGC converte GTP $\rightarrow$ cGMP
    \item cGMP ativa PKG (proteína quinase G)
    \item PKG fosforila alvos
    \item Fosfodiesterase degrada cGMP
\end{enumerate}

\vspace{0.3cm}

\textbf{Características:}
\begin{itemize}
    \item NO é gasoso, lipofilico
    \item Meia-vida curta ($\sim$ segundos)
    \item Sinalização parácrina
\end{itemize}

\column{0.5\textwidth}
\begin{center}
\begin{tikzpicture}[node distance=1.3cm, scale=0.7, transform shape,
    box/.style={rectangle, draw, text width=2cm, text centered, rounded corners, minimum height=0.7cm, font=\small}]

    \node[box, fill=laranjacel!30] (stim) {Estímulo\\(ACh, Shear)};
    \node[box, fill=verdeacido!30, below of=stim] (nos) {NOS\\(NO sintase)};
    \node[box, fill=azulclaro!30, below of=nos] (no) {NO};
    \node[box, fill=roxodna!30, below of=no] (sgc) {sGC\\(guanilil ciclase)};
    \node[box, fill=yellow!30, below of=sgc] (cgmp) {cGMP};
    \node[box, fill=red!20, below of=cgmp] (pkg) {PKG};

    \draw[->, thick] (stim) -- (nos);
    \draw[->, thick] (nos) -- node[right, font=\tiny] {L-Arg} (no);
    \draw[->, thick] (no) -- (sgc);
    \draw[->, thick] (sgc) -- node[right, font=\tiny] {GTP} (cgmp);
    \draw[->, thick] (cgmp) -- (pkg);

    \node[below of=pkg, text width=2.5cm, align=center, font=\tiny] {Relaxamento\\Vascular\\(vasodilatação)};
\end{tikzpicture}
\end{center}

\begin{exampleblock}{Aplicação Clínica}
\textbf{Viagra (sildenafil):} inibe fosfodiesterase que degrada cGMP $\rightarrow$ prolonga efeito vasodilatador!
\end{exampleblock}
\end{columns}
\end{frame}

% ============================================================================
% SEÇÃO 6: DINÂMICA E PROPRIEDADES EMERGENTES
% ============================================================================
\section{Dinâmica e Propriedades Emergentes}

\begin{frame}{9.6 Dinâmica Temporal de Sinalização}
\begin{block}{Padrões Temporais}
\begin{itemize}
    \item \textbf{Transiente:} resposta rápida que retorna ao basal
    \item \textbf{Sustentada:} resposta mantida enquanto estímulo presente
    \item \textbf{Adaptação:} resposta diminui apesar de estímulo constante
    \item \textbf{Oscilações:} ciclos periódicos de ativação/desativação
\end{itemize}
\end{block}

\begin{center}
\begin{tikzpicture}[scale=0.8]
    % Transient
    \begin{scope}
        \draw[->] (0,0) -- (3,0) node[right, font=\tiny] {t};
        \draw[->] (0,0) -- (0,2);
        \draw[thick, azulescuro] (0,0) -- (0.3,0) -- (0.5,1.5) -- (1,1.2) -- (2,0.2) -- (3,0);
        \node[below, font=\small] at (1.5,-0.5) {Transiente};
    \end{scope}

    % Sustained
    \begin{scope}[xshift=4cm]
        \draw[->] (0,0) -- (3,0) node[right, font=\tiny] {t};
        \draw[->] (0,0) -- (0,2);
        \draw[thick, verdeacido] (0,0) -- (0.3,0) -- (0.5,1.5) -- (3,1.5);
        \node[below, font=\small] at (1.5,-0.5) {Sustentada};
    \end{scope}

    % Adaptation
    \begin{scope}[xshift=8cm]
        \draw[->] (0,0) -- (3,0) node[right, font=\tiny] {t};
        \draw[->] (0,0) -- (0,2);
        \draw[thick, laranjacel] (0,0) -- (0.3,0) -- (0.5,1.5) -- (1.5,0.5) -- (3,0.2);
        \node[below, font=\small] at (1.5,-0.5) {Adaptação};
    \end{scope}

    % Oscillations
    \begin{scope}[xshift=12cm]
        \draw[->] (0,0) -- (3,0) node[right, font=\tiny] {t};
        \draw[->] (0,0) -- (0,2);
        \draw[thick, roxodna, domain=0.3:3, samples=50]
            plot (\x, {0.8 + 0.7*sin(360*\x)});
        \node[below, font=\small] at (1.5,-0.5) {Oscilações};
    \end{scope}
\end{tikzpicture}
\end{center}
\end{frame}

\begin{frame}{Feedback Loops: Positivo e Negativo}
\begin{columns}
\column{0.5\textwidth}
\begin{block}{Feedback Negativo}
\begin{center}
\begin{tikzpicture}[scale=0.7]
    \node[draw, circle, fill=azulclaro!30] (a) at (0,0) {A};
    \node[draw, circle, fill=verdeacido!30] (b) at (2.5,0) {B};

    \draw[->, thick] (a) to [out=30, in=150] (b);
    \draw[-|, thick, red] (b) to [out=210, in=330] (a);
\end{tikzpicture}
\end{center}

\textbf{Propriedades:}
\begin{itemize}
    \item Estabilidade
    \item Homeostase
    \item Adaptação
    \item Reduz variabilidade
\end{itemize}

\textbf{Exemplos:}
\begin{itemize}
    \item Inibição por produto
    \item Desensibilização de receptores
    \item I-Smads na via TGF-$\beta$
\end{itemize}
\end{block}

\column{0.5\textwidth}
\begin{block}{Feedback Positivo}
\begin{center}
\begin{tikzpicture}[scale=0.7]
    \node[draw, circle, fill=azulclaro!30] (a) at (0,0) {A};
    \node[draw, circle, fill=verdeacido!30] (b) at (2.5,0) {B};

    \draw[->, thick] (a) to [out=30, in=150] (b);
    \draw[->, thick, verdeacido] (b) to [out=210, in=330] (a);
\end{tikzpicture}
\end{center}

\textbf{Propriedades:}
\begin{itemize}
    \item Amplificação
    \item Biestabilidade
    \item Memória
    \item Decisões irreversíveis
\end{itemize}

\textbf{Exemplos:}
\begin{itemize}
    \item CICR (Ca$^{2+}$)
    \item Ativação de caspases (apoptose)
    \item Ciclo celular (checkpoints)
\end{itemize}
\end{block}
\end{columns}

\vspace{0.3cm}

\importante{
Combinações de feedbacks positivo e negativo geram oscilações!
}
\end{frame}

\begin{frame}{Gradientes Espaciais e Morfógenos}
\definicao{Morfógeno}{
Molécula sinalizadora que forma gradiente de concentração e especifica diferentes destinos celulares dependendo da concentração local.
}

\begin{center}
\begin{tikzpicture}[scale=0.9]
    % Source
    \draw[thick, fill=laranjacel!50] (0,0) rectangle (0.5,3);
    \node[left, font=\small] at (0,1.5) {Fonte};

    % Gradient
    \foreach \x in {1,2,...,8} {
        \pgfmathsetmacro{\opacity}{100 - 10*\x}
        \draw[fill=azulclaro!\opacity] (\x*0.5,0) rectangle (\x*0.5+0.5,3);
    }

    % Concentration profile
    \draw[->, thick] (0,-1) -- (5,-1) node[right, font=\small] {Distância};
    \draw[->, thick] (0,-1) -- (0,-0.5) node[above, font=\small] {[M]};

    \draw[thick, roxodna, domain=0:4.5, samples=50]
        plot (\x, {-1 + 2*exp(-0.8*\x)});

    % Thresholds
    \draw[dashed, red] (0,-0.5) -- (4.5,-0.5);
    \draw[dashed, red] (0,-0.7) -- (4.5,-0.7);

    \node[right, font=\tiny] at (4.5,-0.5) {Threshold 1};
    \node[right, font=\tiny] at (4.5,-0.7) {Threshold 2};

    % Cell fates
    \node[below, font=\small, align=center] at (0.5,-1.8) {Tipo\\Celular A};
    \node[below, font=\small, align=center] at (2,-1.8) {Tipo\\Celular B};
    \node[below, font=\small, align=center] at (3.5,-1.8) {Tipo\\Celular C};
\end{tikzpicture}
\end{center}

\exemplo{Morfógenos em Desenvolvimento}{
\begin{itemize}
    \item \textbf{Bicoid:} padrão ântero-posterior em \textit{Drosophila}
    \item \textbf{Shh (Sonic Hedgehog):} tubo neural e membros
    \item \textbf{BMP:} padrão dorso-ventral
\end{itemize}
}
\end{frame}

\begin{frame}{Oscilações em Sistemas de Sinalização}
\begin{columns}
\column{0.5\textwidth}
\textbf{Requisitos para Oscilações:}
\begin{itemize}
    \item Feedback negativo com atraso
    \item Não-linearidade (Hill, saturação)
    \item Separação de escalas temporais
\end{itemize}

\vspace{0.3cm}

\textbf{Exemplos Biológicos:}
\begin{enumerate}
    \item \textbf{NF-$\kappa$B:} oscilações em resposta a TNF-$\alpha$
    \item \textbf{p53:} pulsos em resposta a dano ao DNA
    \item \textbf{Ca}$^{2+}$: oscilações citosólicas
    \item \textbf{ERK:} pulsos vs sustentado
    \item \textbf{Hes1:} oscilações no relógio de segmentação
\end{enumerate}

\column{0.5\textwidth}
\begin{exampleblock}{Caso: p53}
\begin{center}
\begin{tikzpicture}[scale=0.7]
    \node[draw, circle, fill=azulclaro!30] (p53) at (0,0) {p53};
    \node[draw, circle, fill=verdeacido!30] (mdm2) at (3,0) {Mdm2};

    \draw[->, thick] (p53) to [out=30, in=150] node[above, font=\tiny] {ativa} (mdm2);
    \draw[-|, thick, red] (mdm2) to [out=210, in=330] node[below, font=\tiny] {degrada} (p53);
\end{tikzpicture}
\end{center}

\vspace{0.2cm}

p53 ativa Mdm2 $\rightarrow$ Mdm2 degrada p53 $\rightarrow$ oscilações!

\vspace{0.3cm}

\textbf{Significado Funcional:}
\begin{itemize}
    \item Pulsos de p53 permitem reparo do DNA
    \item p53 sustentado $\rightarrow$ apoptose
    \item Número de pulsos codifica severidade do dano
\end{itemize}
\end{exampleblock}
\end{columns}
\end{frame}

\begin{frame}{Adaptação e Desensibilização}
\begin{block}{Mecanismos de Adaptação}
\begin{enumerate}
    \item \textbf{Fosforilação de receptores:} reduz afinidade/atividade
    \item \textbf{Internalização de receptores:} endocitose mediada
    \item \textbf{Indução de fosfatases/inibidores:} feedback negativo
    \item \textbf{Depleção de segundos mensageiros:} exaustão de estoques
\end{enumerate}
\end{block}

\vspace{0.3cm}

\exemplo{Desensibilização de GPCRs}{
\begin{enumerate}
    \item GPCR ativo é fosforilado por \textbf{GRKs} (GPCR kinases)
    \item \textbf{$\beta$-arrestina} liga ao receptor fosforilado
    \item Bloqueia acoplamento com proteína G
    \item Promove internalização via endocitose
    \item Receptor pode ser reciclado ou degradado
\end{enumerate}
}

\vspace{0.3cm}

\importante{
Adaptação permite que células detectem \textbf{mudanças} no estímulo, não apenas intensidade absoluta!
}
\end{frame}

\begin{frame}{Robustez e Sensibilidade}
\begin{columns}
\column{0.5\textwidth}
\begin{block}{Robustez}
Capacidade do sistema manter função apesar de perturbações.

\vspace{0.2cm}

\textbf{Mecanismos:}
\begin{itemize}
    \item Feedback negativo
    \item Redundância (vias paralelas)
    \item Buffering (sequestro)
    \item Compartimentalização
\end{itemize}
\end{block}

\begin{exampleblock}{Exemplo}
Via MAPK é robusta a variações em concentração de componentes devido a feedbacks múltiplos.
\end{exampleblock}

\column{0.5\textwidth}
\begin{block}{Sensibilidade}
Capacidade de responder a pequenas mudanças no estímulo.

\vspace{0.2cm}

\textbf{Mecanismos:}
\begin{itemize}
    \item Amplificação (cascatas)
    \item Ultrassensibilidade (Hill)
    \item Feedback positivo
    \item Cooperatividade
\end{itemize}
\end{block}

\begin{exampleblock}{Exemplo}
Cascata MAPK amplifica sinal $>$100$\times$ e exibe ultrassensibilidade (Hill $n \sim 5$).
\end{exampleblock}
\end{columns}

\vspace{0.3cm}

\importante{
Sistemas de sinalização equilibram \textbf{robustez} (resistência a ruído) com \textbf{sensibilidade} (detecção de sinais)!
}
\end{frame}

% ============================================================================
% SEÇÃO 7: EXERCÍCIOS
% ============================================================================
\section{Exercícios}

\begin{frame}{Exercícios - Parte 1}
\begin{block}{Questão 1: Ligação Receptor-Ligante}
Um receptor tem $K_d = 5$ nM e $B_{max} = 1000$ receptores/célula.
\begin{enumerate}
    \item Calcule a ocupação quando [Ligante] = 5 nM
    \item Que concentração de ligante resulta em 90\% de ocupação?
    \item Se $k_{on} = 10^6$ M$^{-1}$s$^{-1}$, qual é $k_{off}$?
\end{enumerate}
\end{block}

\begin{block}{Questão 2: Amplificação em Cascata}
Uma cascata MAPK tem 3 níveis. Cada quinase ativada fosforila 10 moléculas do próximo nível. Se 5 moléculas de MAPKKK são ativadas:
\begin{enumerate}
    \item Quantas moléculas de MAPKK são ativadas?
    \item Quantas moléculas de MAPK são ativadas?
    \item Qual é o fator de amplificação total?
\end{enumerate}
\end{block}
\end{frame}

\begin{frame}{Exercícios - Parte 2}
\begin{block}{Questão 3: Ultrassensibilidade}
Uma quinase exibe comportamento ultrassensível com coeficiente de Hill $n_H = 4$ e $K_{0.5} = 10$ $\mu$M.
\begin{enumerate}
    \item Escreva a equação de Hill para a ativação
    \item Calcule a fração ativa quando [S] = 5, 10, 15 $\mu$M
    \item Compare com Michaelis-Menten ($n_H = 1$)
\end{enumerate}
\end{block}

\begin{block}{Questão 4: Segundos Mensageiros}
cAMP é produzido a $v_{synth} = 0.5$ $\mu$M/s e degradado com constante $k_{deg} = 0.1$ s$^{-1}$.
\begin{enumerate}
    \item Escreva a equação diferencial para [cAMP]
    \item Calcule a concentração no estado estacionário
    \item Quanto tempo leva para atingir 90\% do steady-state após estimulação?
\end{enumerate}
\end{block}
\end{frame}

\begin{frame}{Exercícios - Parte 3}
\begin{alertblock}{Questão 5: Simulação de Cascata MAPK}
Use Python para simular a cascata MAPK:
\begin{enumerate}
    \item Implemente o modelo de 3 níveis (MAPKKK $\rightarrow$ MAPKK $\rightarrow$ MAPK)
    \item Inclua fosforilação e desfosforilação
    \item Simule a dinâmica temporal para um estímulo em degrau
    \item Gere curva dose-resposta variando intensidade do estímulo
    \item Calcule o coeficiente de Hill efetivo a partir da curva dose-resposta
\end{enumerate}
\end{alertblock}

\begin{block}{Questão 6: Oscilações de Ca$^{2+}$}
Modifique o código de oscilações de Ca$^{2+}$ para:
\begin{enumerate}
    \item Incluir feedback negativo via inativação do receptor IP$_3$
    \item Simular ondas de Ca$^{2+}$ (1D ou 2D com difusão)
    \item Investigar como parâmetros afetam frequência de oscilação
\end{enumerate}
\end{block}
\end{frame}

\begin{frame}{Exercícios - Parte 4}
\begin{alertblock}{Projeto Prático: Modelagem de Via Completa}
Escolha uma via de sinalização (MAPK, PI3K/Akt, JAK-STAT) e:
\begin{enumerate}
    \item Desenhe diagrama completo da via com todas as interações
    \item Identifique pontos de regulação (feedbacks, crosstalk)
    \item Construa modelo matemático (ODEs) com 5-10 componentes
    \item Implemente o modelo em Python usando \texttt{scipy.integrate.odeint}
    \item Simule perturbações (knockout de componente, inibidor farmacológico)
    \item Analise sensibilidade dos parâmetros
    \item Compare previsões do modelo com dados experimentais da literatura
    \item Gere visualizações: dinâmica temporal, espaço de fases, bifurcações
\end{enumerate}
\end{alertblock}
\end{frame}

% ============================================================================
% SEÇÃO 8: REFERÊNCIAS
% ============================================================================
\section{Referências}

\begin{frame}{Referências Principais}
\begin{thebibliography}{99}
\bibitem{alberts} Alberts B. et al. (2022). \textit{Molecular Biology of the Cell}, 7th ed. Garland Science.

\bibitem{lodish} Lodish H. et al. (2021). \textit{Molecular Cell Biology}, 9th ed. W.H. Freeman.

\bibitem{alon} Alon U. (2019). \textit{An Introduction to Systems Biology: Design Principles of Biological Circuits}, 2nd ed. CRC Press.

\bibitem{kholodenko} Kholodenko B.N. (2006). Cell-signalling dynamics in time and space. \textit{Nature Reviews Molecular Cell Biology} 7:165-176.

\bibitem{ferrell} Ferrell J.E. (1996). Tripping the switch fantastic: how a protein kinase cascade can convert graded inputs into switch-like outputs. \textit{Trends in Biochemical Sciences} 21:460-466.
\end{thebibliography}
\end{frame}

\begin{frame}{Leitura Complementar}
\begin{block}{Livros}
\begin{itemize}
    \item Tyson J.J., Novak B. (2020). \textit{The Cell Cycle: Principles of Control}. Oxford.
    \item Klipp E. et al. (2016). \textit{Systems Biology: A Textbook}, 2nd ed.
    \item Fall C.P. et al. (2002). \textit{Computational Cell Biology}. Springer.
\end{itemize}
\end{block}

\begin{block}{Artigos Clássicos}
\begin{itemize}
    \item Goldbeter A., Koshland D.E. (1981). Ultrasensitivity in biochemical systems. \textit{PNAS}.
    \item Bhalla U.S., Iyengar R. (1999). Emergent properties of networks of signaling pathways. \textit{Science}.
    \item Dolmetsch R.E. et al. (1998). Calcium oscillations increase efficiency of gene expression. \textit{Nature}.
\end{itemize}
\end{block}

\begin{block}{Recursos Online}
\begin{itemize}
    \item \textbf{KEGG Pathway:} \url{https://www.genome.jp/kegg/pathway.html}
    \item \textbf{Reactome:} \url{https://reactome.org}
    \item \textbf{SignaLink:} \url{http://signalink.org}
    \item \textbf{PathwayCommons:} \url{https://www.pathwaycommons.org}
\end{itemize}
\end{block}
\end{frame}

% ============================================================================
% SLIDE FINAL
% ============================================================================
\begin{frame}
\begin{center}
{\Huge Obrigado!}

\vspace{1cm}

{\Large Capítulo 9: Signaling Systems}

\vspace{1cm}

\textbf{Próximo Capítulo:}\\
Metabolic Systems (Sistemas Metabólicos)

\vspace{1cm}

\textit{Dúvidas? Perguntas?}
\end{center}
\end{frame}

\end{document}
